%% en/sections/07_public_history.tex — Public History (full EN)

\chapter{Public History in the Classroom}
\markboth{Public History}{Public History}

\lettrine[lines=3]{P}{laying} the Sicilian Vespers in class is not
``using games to teach history.'' It is something more ambitious:
placing students in the position of people who must make historical
choices without knowing how things will end. This changes their
relationship with the past --- and with the present.

This section is addressed to teachers and GM-educators who wish to
use \textit{Vespri 1282} in a school context. It is primarily aimed
at \textbf{upper secondary school} (ages 14--19: general, scientific,
humanities, and technical programmes), though most guidance is adaptable
to other levels.

\nota[An Honest Premise]{
  Roleplaying in the classroom works when the teacher is willing to
  cede some control over the experience. You do not know exactly what
  your students will say, what they will decide, how they will react.
  That is the point. If you want to control every variable, this tool
  is not for you. If you want to be surprised and to surprise, read on.
}

\secsep

\section{Why the Sicilian Vespers in Class}

\subsection{Curricular Connections}

\textit{Vespri 1282} connects directly to core topics present in upper
secondary history curricula across Europe:

\begin{itemize}
  \item \textbf{Medieval history}: the Mediterranean as a space of conflict
    and contact; the Papacy and temporal power; 13th-century European
    monarchies
  \item \textbf{Medieval geopolitics}: competition between Anjou, Aragon,
    and Byzantium; the crusade as political instrument
  \item \textbf{History of southern Italy and Sicily}: from the Arab emirate
    to the Normans, Hohenstaufen, and Angevins; Sicily as a multicultural space
  \item \textbf{Civic education}: power and resistance; rights and oppression;
    collective violence as political response
  \item \textbf{Historiography}: the debate between event-history and structural
    history; primary sources and their limits
\end{itemize}

\subsection{Transferable Skills}

Beyond disciplinary content, \textit{Vespri 1282} develops:

\begin{itemize}
  \item \textbf{Critical thinking}: evaluating partial and contradictory information
  \item \textbf{Historical empathy}: adopting the perspective of people in contexts
    very different from one's own
  \item \textbf{Cooperation}: working with a group toward shared goals,
    managing internal conflicts
  \item \textbf{Communication}: arguing, persuading, negotiating in real time
  \item \textbf{Counterfactual reasoning}: imagining plausible historical
    alternatives and evaluating their consequences
\end{itemize}

\secsep

\section{Session Structure}

\subsection{Recommended Timing}

\begin{tabular}{lll}
\toprule
\cinzel Phase & \cinzel Duration & \cinzel When \\
\midrule
Historical preparation & 2--3 lessons & weeks before \\
Rules introduction & 1 lesson & day before or same morning \\
One-shot & 3--4 hours & afternoon block or full day \\
Debrief & 1--2 hours & immediately after \\
Written task & homework & following week \\
\bottomrule
\end{tabular}

\subsection{Historical Preparation}

Before the session, students need a basic knowledge of the context.
Recommended:

\begin{enumerate}
  \item \textbf{Guided reading} of the Historical Context chapter of this
    manual, in class or at home
  \item \textbf{Primary source analysis}: a passage from Giovanni Villani
    or Bartholomew of Neocastro on the causes of the uprising
  \item \textbf{Class discussion}: \textit{Who was clearly wrong? Who was
    partially right? Who had no choice?}
  \item \textbf{Role assignment}: choose or assign factions before the session,
    so students know who they ``are'' and can prepare
\end{enumerate}

\nota[On Role Assignment]{
  For large classes (30+ students), divide into four groups and assign
  factions. If possible, let each group choose their own playbooks
  internally. The choice process is already an activity:
  \textit{who do you want to be in this story?}
}

\subsection{Session Day}

\textbf{Morning (if a full day):}
\begin{itemize}[nosep]
  \item 9:00--10:00 --- Historical review, rules introduction, questions
  \item 10:00--10:15 --- Break; GMs position at their tables
  \item 10:15--10:30 --- In-world introduction: each GM reads their
    faction's opening text aloud
\end{itemize}

\textbf{Play session:}
\begin{itemize}[nosep]
  \item 10:30--12:00 --- First act (through the Rumour)
  \item 12:00--13:00 --- Lunch
  \item 13:00--14:30 --- Second act (through the Crisis)
  \item 14:30--15:30 --- Third act (through the Bells)
  \item 15:30--16:00 --- Resolution and outcome reveal
\end{itemize}

\secsep

\section{Opening Texts by Faction}
\textit{(To be read aloud by the GM at session start)}

\subsection{Opening: The Conspirators}

\begin{citbox}
\textit{Palermo, February 1282. The night is cold. You are here ---
in this place, at this moment --- because someone trusted you enough
to tell you the truth.}

\textit{The truth is this: in a few weeks, if everything goes as it
should, Palermo will no longer be Angevin. The truth is also this:
not everything is going as it should. And the most uncomfortable truth
is that none of you knows exactly what will happen --- only what
\textbf{you} will do.}

\textit{You chose to be here. Now you will choose the rest.}
\end{citbox}

\subsection{Opening: The Barons}

\begin{citbox}
\textit{Palermo, February 1282. You received a message that should not
have been sent to you. Or perhaps it should --- perhaps someone decided
it was time for you to know.}

\textit{The proposal is simple: something is moving. Something large.
You can be part of it --- or you can stand aside and hope it does not
sweep you away. But the message also says one more thing: the time to
decide is now. Not tomorrow. Not next week.}

\textit{Now.}
\end{citbox}

\subsection{Opening: The Clergy}

\begin{citbox}
\textit{Palermo, February 1282. The bells have just rung Compline.
The city has gone quiet. But this quiet is not peace --- it is the
quiet of someone holding their breath.}

\textit{You serve God. You serve Rome. You serve your people. Tonight,
for the first time, you are not certain these three things point in
the same direction. Tonight you may do something that cannot be undone.}

\textit{Pray. And then decide.}
\end{citbox}

\subsection{Opening: The People}

\begin{citbox}
\textit{Palermo, February 1282. It is morning. The market is open.
The French soldiers are already in the square --- as always, for
sixteen years, as always.}

\textit{You know nothing of any conspiracy. You only know what you
know: the taxes have come again. Your cousin was beaten last week for
no reason. The price of bread has gone up again. And someone told you,
yesterday evening, that in a few weeks everything might change.}

\textit{You do not know whether you believe it. But you were asked to
be here. And here you are.}
\end{citbox}

\secsep

\section{Structured Debrief}

The debrief is the most important part of the educational activity.
Without a debrief, the game is only an experience. With a debrief,
it becomes an opportunity for deep learning.

\subsection{Phase 1: De-roll (10--15 minutes)}

Before discussing history or meanings, it is essential to help students
exit their characters.

\textbf{De-roll questions:}
\begin{itemize}[nosep]
  \item ``What was your character feeling in the last scene?''
  \item ``What were \textbf{you} feeling in that same scene?''
  \item ``Was there a choice you made that you are proud of?
    One you regret?''
\end{itemize}

This step is especially important if the session produced emotionally
intense moments. Do not skip it.

\subsection{Phase 2: Historical Reflection (30--45 minutes)}

\textbf{Discussion questions:}

\begin{enumerate}
  \item \textit{Did your outcome match the historical one? If not,
    what went differently? Why?}

  \item \textit{Which historical figures did you encounter? How did
    you portray them? Did they match what you knew about them?}

  \item \textit{Was there a moment when you understood something about
    the historical situation that you had not understood before? What?}

  \item \textit{Was the Vespers uprising ``just''? Was it ``inevitable''?
    These are two different questions --- let us discuss them separately.}

  \item \textit{What could Charles of Anjou have done to prevent the
    uprising? Was it possible?}
\end{enumerate}

\subsection{Phase 3: Connections to the Present (15--20 minutes)}

\begin{enumerate}
  \item \textit{Are there situations in the contemporary world that
    recall the power structure you played? (Caution: do not look for
    direct analogies --- look for similar structures)}

  \item \textit{Who writes the history of the Vespers? The sources you
    read --- Villani, Bartholomew of Neocastro --- who were they?
    Why were they writing?}

  \item \textit{The People won. Yet their names are not in the chronicles.
    How is this possible? Is it just?}
\end{enumerate}

\secsep

\section{Written Tasks}

\subsection{Task Options (student's choice)}

\textbf{A --- The Character's Diary}\\
Write your character's diary for the three weeks before 30 March 1282.
Include their motivations, fears, the choices made during the session,
and what they think will happen next.

\textbf{Requirements:} 600--800 words. Historical accuracy (no
anachronisms). At least one reference to a real historical source
(cited in bibliography).

\filetto

\textbf{B --- The Historiographical Analysis}\\
Compare two interpretations of the ``Aragonese conspiracy'': Runciman's
(organised plot) and Abulafia's (spontaneous revolt). Which seems more
convincing, based on what you have studied and what you experienced in
the game? Argue your position.

\textbf{Requirements:} 800--1000 words. Argumentative structure (thesis,
arguments, counter-arguments, conclusion). At least two bibliographical
sources.

\filetto

\textbf{C --- The Alternative Source}\\
Imagine you are a 14th-century historian writing your own chronicle of
the Sicilian Vespers. Your perspective is different from Villani's: you
are a Sicilian friar, a Genoese merchant, or an Angevin official who
escaped the night of 30 March. Write your chronicle.

\textbf{Requirements:} 700--900 words. Style coherent with medieval
writing (first person, religious references, formal but not academic
language). A brief author's note explaining the choices made.

\secsep

\section{Teacher Notes}

\subsection{Emotional Management}

The game can produce emotionally intense moments --- especially at the
People table, where characters are vulnerable and consequences are
concrete.

\begin{itemize}
  \item Establish a \textbf{safety word} before the session: a word
    anyone can say to pause the game without explanation. Honour it
    without exception.
  \item The GM at each table should monitor players' emotional state.
    If someone is clearly uncomfortable, offer them the option to take
    a break.
  \item Violence in the game must remain \textbf{narrative} --- never
    physical reenactment. Conflicts are resolved with dice, not bodies.
\end{itemize}

\subsection{Adaptations for Different Classes}

\textbf{Students with learning differences (SEN/SLD):}
\begin{itemize}[nosep]
  \item Provide character sheets with simplified text
  \item Allow manual consultation during play
  \item Assign roles with more dialogic platforms (e.g. Notary, Canon)
    to students who prefer not to be at the centre of attention
\end{itemize}

\textbf{Large classes (35+ students):}
\begin{itemize}[nosep]
  \item Add a fifth ``observer table'' that follows one faction,
    takes notes, and reports during the debrief
  \item Reduce each faction to 5--6 characters, removing the less
    central playbooks
\end{itemize}

\textbf{Reduced time (2 hours instead of 4):}
\begin{itemize}[nosep]
  \item Remove the third act (The Bells) and reveal the outcome based
    on scores after the Crisis only
  \item Focus each faction on a single main objective
\end{itemize}

\filettodoppio
